%
% drehung.tex -- Drehung für die Berechnung des Integrals
%
% (c) 2021 Prof Dr Andreas Müller, OST Ostschweizer Fachhochschule
%
\documentclass[tikz]{standalone}
\usepackage{amsmath}
\usepackage{times}
\usepackage{txfonts}
\usepackage{pgfplots}
\usepackage{csvsimple}
\usetikzlibrary{arrows,intersections,math,calc,patterns,angles}
\usetikzlibrary{arrows.meta,positioning,backgrounds,quotes}
\definecolor{darkred}{rgb}{0.8,0,0}
\definecolor{darkgreen}{rgb}{0,0.6,0}
\begin{document}
\def\skala{1}
\begin{tikzpicture}[>=latex,thick,scale=\skala]

    \colorlet{Ecol}{orange!90!black}
    \colorlet{Icol}{blue!50!black}
    \colorlet{Rcol}{green!50!black}

    \tikzstyle{vector}=[->,very thick,line cap=round]

    \def\xmin{1.8}
    \def\xmax{3}
    \def\ymax{2.5}
    \def\R{2.7} 
    \def\ang{45}

    \coordinate (O) at (0,0);
    \coordinate (X) at (1.4*\xmax,0);
    \coordinate (Y) at (0,1.3*\ymax);
    \coordinate (R) at (\ang:\R);
    \coordinate (Rx) at (\R,0);
    
    % AXIS
    \draw[->] (-0.3*\xmax,0) -- (X) node[below] {$\text{Re}$};
    \draw[->] (0,-0.3*\ymax) -- (Y) node[left] {$\text{Im}$};
    
    % PHASORS
    \draw[vector,Ecol] (O) -- (R) %node[above right=-0.2] {$\VR$};
        node[midway, above, sloped] {$\text{gedreht}$};

    \draw pic[draw, ->, Icol, very thick, angle radius = 2.7cm] 
        {angle=Rx--O--R};
    \draw[thick, Icol] (2.5,0.9) node[above right] {$\text{Bogen}$};

    \draw[fill=red] (Rx) circle (2pt);
    \draw[Rcol] (Rx) node[below right] {$R$};

    \draw[vector,Rcol] (O) -- (Rx)
        node[midway, below] {$\text{horizontal}$};

    \draw pic[draw, ->, angle radius = 0.8cm] {angle=Rx--O--R};
    \draw (0.3,0) node[above right] {$\vartheta$};

\end{tikzpicture}
\end{document}

